%!TEX program = xelatex
\def\ano{2021}
\def\semestre{1}
\def\disciplina{\'Algebra 1}
\def\turma{C}
\def\numerosemana{11}

\documentclass[12pt]{exam}

\usepackage{caption}
\usepackage{amssymb}
\usepackage{amsmath,amsfonts,amsthm,amstext}
\usepackage[brazil]{babel}
% \usepackage[latin1]{inputenc}
\usepackage{graphicx}
\graphicspath{{/home/jfreitas/GitHub_Repos/Algebra-1/Pictures/}{D:/OneDrive - unb.br/imagens-latex/}}
\usepackage{enumitem}
\usepackage{multicol}
\usepackage{answers}
\usepackage{tikz,ifthen}
\usetikzlibrary{lindenmayersystems}
\usetikzlibrary[shadings]
\Newassociation{solucao}{Solution}{ans}
\newtheorem{exercicio}{}

\setlength{\topmargin}{-1.0in}
\setlength{\oddsidemargin}{0in}
\setlength{\textheight}{10.1in}
\setlength{\textwidth}{6.5in}
\setlength{\baselineskip}{12mm}

\extraheadheight{0.7in}
\firstpageheadrule
\runningheadrule
\lhead{
        \begin{minipage}[c]{1.7cm}
        \includegraphics[width=1.7cm]{unb.pdf}
        \end{minipage}%
        \hspace{0pt}
        \begin{minipage}[c]{4in}
          {Universidade de Bras{\'\i}lia} --
          {Departamento de Matem{\'a}tica}
\end{minipage}
\vspace*{-0.8cm}
}
% \chead{Universidade de Bras{\'\i}lia - Departamento de Matem{\'a}tica}
% \rhead{}
% \vspace*{-2cm}

\extrafootheight{.5in}
\footrule
\lfoot{\disciplina\ - \semestre$^o$/\ano}
\cfoot{}
\rfoot{P\'agina \thepage\ de \numpages}

\newcounter{exercicios}
\renewcommand{\theexercicios}{\arabic{exercicios}}

\newenvironment{questao}[1]{
\refstepcounter{exercicios}
\ifx&#1&
\else
   \label{#1}
\fi
\noindent\textbf{Exerc{\'\i}cio {\theexercicios}:}
}

\newcommand{\resp}[1]{
\noindent{\bf Exerc{\'\i}cio #1: }}
\newcommand{\nsub}{\varsubsetneq}
\newcommand{\vaz}{\emptyset}
\newcommand{\im}{{\rm Im\,}}
\newcommand{\sub}{\subseteq}
\newcommand{\n}{\mathbb{N}}
\newcommand{\z}{\mathbb{Z}}
\newcommand{\rac}{\mathbb{Q}}
\newcommand{\real}{\mathbb{R}}
\newcommand{\complex}{\mathbb{C}}
\newcommand{\cp}[1]{\mathbb{#1}}
\newcommand{\ch}{\mbox{\textrm{car\,}}\nobreak}
\newcommand{\vesp}[1]{\vspace{ #1  cm}}
\newcommand{\compcent}[1]{\vcenter{\hbox{$#1\circ$}}}
\newcommand{\comp}{\mathbin{\mathchoice
{\compcent\scriptstyle}{\compcent\scriptstyle}
{\compcent\scriptscriptstyle}{\compcent\scriptscriptstyle}}}

\begin{document}

\begin{center}

    {\Large\bf \disciplina\ - Turma \turma\ -- \semestre$^{o}$/\ano} \\ \vspace{9pt} {\large\bf
        Lista de Exerc{\'\i}cios -- Semana \numerosemana}\\ \vspace{9pt} Prof. Jos{\'e} Ant{\^o}nio O. Freitas
    \end{center}
    \hrule

    \vspace{.6cm}

    \questao{} Verifique se s\~ao subgrupos:
    \begin{enumerate}[label=({\alph*})]
      \item $H = \{x \in \rac \mid x > 0\}$ de $(\rac^*,\cdot)$.

      \item $H = \left\{\dfrac{1 + 2m}{1 + 2n} \mid m, n \in \z\right\}$ de $(\rac^*,\cdot)$.

      \item $H = \{\cos\theta + i\sin\theta \mid \theta \in \rac\}$ de $(\complex^*,\cdot)$.

      \item $H = \{0, \pm 2, \pm 4, \pm 6, \dots\}$ de $(\z,+)$.

      \item $H = \{0, \pm 2, \pm 4, \pm 6, \dots\}$ do grupo $(\rac - \{1\},\star)$ onde $\star$ \'e definida como $x \star y = x + y - xy$.

      \item $H = \{a + b\sqrt{2} \mid a, b \in \rac\}$ de $(\real,+)$.

      \item $H = \{a + b\sqrt{2} \in \real^* \mid a, b \in \rac\}$ de $(\real^*,\cdot)$.

      \item $H = \{a + b\sqrt[3]{2} \mid a, b \in \rac\}$ de $(\real,+)$.

      \item $H = \{a + b\sqrt[3]{2} \in \real^* \mid a, b \in \rac\}$ de $(\real^*,\cdot)$.
    \end{enumerate}

    \vspace{.3cm}

    \questao{} Determine todos os subgrupos do grupo aditivo $\z_4$.

    \vspace{.3cm}

    \questao{} Determine todos os subgrupos de $S_3$.

    \vspace{.3cm}

    \questao{} Seja
    \[
      GL_2(\real) = \left\{ \begin{pmatrix}
          x & y\\z & t
      \end{pmatrix} \mid x, y, z, t \in \real,\ \det(A) \ne 0\right\}.
    \]
    \begin{enumerate}[label=({\alph*})]
      \item Mostre que $GL_2(\real)$ com a opera\c{c}\~ao de multiplica\c{c}\~ao de matrizes \'e um grupo. Esse grupo \'e abeliano?

      \item Seja
      \[
          H = \left\{ \begin{pmatrix}
              \cos a & \sin a\\ - \sin a & \cos a
          \end{pmatrix} \mid a \in \real\right\}.
      \]
      Mostre que $H$ \'e um subgrupo de $GL_2(\real)$.

      \item Seja
      \[
          K = \left\{ \begin{pmatrix}
              a & b\\ -b & a
          \end{pmatrix} \mid a, b \in \real \mbox{ e n\~ao nulos simultaneamente}\right\}.
      \]
      Mostre que $K$ \'e um subgrupo de $GL_2(\real)$.
    \end{enumerate}

    \questao{} Sejam $H$ e $K$ subgrupos de um grupo $G$ (com nota{\c c}{\~a}o
    multiplicativa).
    \begin{enumerate}[label=({\alph*})]
      \item Mostre que $H\cap K$ tamb{\'e}m {\'e} subgrupo de $G$.

      \item Seja $g\in G$ um elemento fixado. Mostre que o conjunto
      $g^{-1}Hg=\{ g^{-1}xg \mid x\in H \} $ {\'e} um subgrupo de $G$.

      \item Prove que $H\cup K$ {\'e} subgrupo de $G$ se, e somente se,
      $H\subseteq K$ ou $K\subseteq H$.

      \item Demonstre que $HK=\{hk \mid h\in H, k\in K\}$ {\'e} subgrupo
      de $G$ se, e somente se, $HK=KH$.

      [\emph{Nota: $HK=KH$ \textbf{n{\~a}o} quer dizer que $hk=kh$,
      para todo $h\in H, k\in K$; significa que $hk=k_1h_1 \in KH$ e $kh=h_2k_2 \in
      HK$, para todo $h\in H, k\in K$.}]
    \end{enumerate}

    \vspace{.3cm}
    \questao{} Seja $G$ um grupo com nota\c{c}\~ao multiplicativa e $a$ um elemento de $G$. Prove que $N(a) = \{x \in G \mid ax = xa\}$ \'e um subgrupo de $G$.

    \vspace{.3cm}

    \questao{} Seja $G$ um grupo com nota\c{c}\~ao multiplicativa. Considere o subconjunto $Z(G) = \{x \in G \mid xh = hx, \mbox{ para todo } h \in G\}$. Mostre que:
    \begin{enumerate}[label=({\alph*})]
      \item $Z(G)$ \'e um subgrupo de $G$.

      \item $G$ \'e abeliano se, e somente se, $Z(G) = G$.
    \end{enumerate}

    \vspace{.3cm}

    \questao{} Seja $G$ um grupo. Dado $H \subset G$ um subgrupo defina
    \begin{align}\label{relacao_equivalencia_subgrupo}
        x \cong y \mbox{ se, e somente se, } xy^{-1} \in H
    \end{align}
    para todos $x$, $y \in G$.
    \begin{enumerate}[label={\alph*})]
        \item Mostre que a rela\c{c}\~ao definida em \eqref{relacao_equivalencia_subgrupo} \'e uma rela\c{c}\~ao de equival\^encia.

        \item Se $y \in G$, ent\~ao a classe de equival\^encia determinada por $y$ \'e o conjunto
        \begin{align*}\label{classe_equivalencia_subgrupo}
            Hy = \{ty \mid t \in H\}.
        \end{align*}
    \end{enumerate}

    \vspace{.3cm}

    \questao{} Seja $H$ um subgrupo de $G$. Denote por
    \begin{align*}
        \mathcal{P} &= \{aH \mid a \in G\}\\
        \mathcal{Q} &= \{Hb \mid b \in G\}.
    \end{align*}
    Mostre que a fun\c{c}\~ao $f : \mathcal{P} \to \mathcal{Q}$ dada por $f(aH) = Ha^{-1}$ \'e uma bije\c{c}\~ao.

    \vspace{.3cm}

    \questao{} Determine todas as classes laterais de $H = \{\overline{0}, \overline{3}, \overline{6}, \overline{9}\}$ no grupo aditivo $\z_{12}$.

    \vspace{.3cm}

    \questao{} Construa os seguintes subgrupos:
    \begin{enumerate}[label={\alph*})]
        \item $[-1]$ em $(\rac, +)$

        \item $[6]$ em $(\z, +)$

        \item $[6]$ em $(\rac^*, \cdot)$

        \item $[i]$ em $(\complex, +)$
    \end{enumerate}

    \vspace{.3cm}

    \questao{} Determine todas as classes laterais de $4\z$ no grupo aditivo $\z$.

    \vspace{.3cm}

    \questao{} Seja $S_3$ o grupo de permuta\c{c}\~oes de $A = \{1, 2, 3\}$. Determine todas as classes laterais do subgrupo $H = \{f, g\}$, onde
    \begin{align*}
        f &= \begin{pmatrix}
            1 & 2 & 3\\
            1 & 2 & 3
        \end{pmatrix}\\
        g &= \begin{pmatrix}
            1 & 2 & 3\\
            2 & 1 & 3
        \end{pmatrix}.
    \end{align*}

    \vspace{.3cm}

    \questao{} Considerando $\z$ como subgrupo do grupo aditivo $\rac$, descreva as classes $\z + (-1)$ e $\z + \dfrac{1}{2}$.

    \vspace{.3cm}

    \questao{} \begin{enumerate}[label={\alph*})]
        \item Mostre que
        \[
            H = \left\{\begin{pmatrix}
                c & 0\\
                0 & c
            \end{pmatrix} \mid c \in \real^*\right\}
        \]
        \'e um subgrupo do grupo $GL_2(\real)$.

        \item Mostre que existem infinitas classes de $H$ em $GL_2(\real)$.
    \end{enumerate}

    \vspace{.3cm}

    \questao{} Mostre que os elementos n\~ao nulos de $\z_{13}$ formam um grupo multiplicativo c{\'\i}clico.

    \vspace{.3cm}

    \questao{} Seja $a \ne e$ um elemento do grupo $G$. Mostre que $o(a) = 2$ se, e somente se, $a = a^{-1}$.

    \vspace{.3cm}

    \questao{} Seja $G$ um grupo multiplicativo $x \in G$. Mostre que, se existe um inteiro $n \ge 1$ tal que $x^n = e$, ent\~ao existe um inteiro $m \ge 1$ tal que $x^{-1} = x^m$.

    \vspace{.3cm}

    \questao{} Seja $G$ um grupo multiplicativo e suponha que $a \in G$. Mostre que $o(a) = o(a^{-1}) = o(xax^{-1})$.

    \vspace{.3cm}

    \questao{} Seja $G$ um grupo finito e denote por $e$ seu elemento neutro. Se $x \in G$, mostre que existe $n \in \z$ tal que $x^n = e$.

    \vspace{.3cm}

    \questao{} Sejam $a$ e $b$ elementos de um grupo $G$ finito. Se $[a] = [b]$, prove que $a$ e $b$ t\^em a mesma ordem.

    \vspace{.3cm}

    \questao{} Considere as seguintes matrizes como elementos do grupo $GL_2(\real)$:
    \[
        A = \begin{bmatrix}
            0 & 1\\
            -1 & -1
        \end{bmatrix}\quad
        B = \begin{bmatrix}
            0 & -1\\
            1 & 0
        \end{bmatrix}.
    \]
    Mostre que $A$ tem ordem 3 e que $B$ tem ordem 4. Qual é a ordem de $AB$?

    \vspace{.3cm}

    \questao{} Dê exemplos do seguinte:
    \begin{enumerate}[label={\alph*})]
        \item Um grupo infinito $G$ com um subgrupo $H \ne G$ e tal que a quantidade de classes laterais de $H$ em $G$ seja finita.

        \item Um grupo infinito $G$ com um subgrupo $H \ne \{e\}$ tal que $|H|$ é finita.
    \end{enumerate}

    \vspace{.3cm}

    \questao{} Seja $G$ um grupo abeliano. Mostre que o conjunto dos elementos de $G$ que possuem ordem finita é um subgrupo de $G$.

    \vspace{.3cm}

    \questao{} Seja $G$ um grupo finito de ordem $n$ tal que todo elemento diferente da identidade têm ordem 2.
    \begin{enumerate}[label={\alph*})]
        \item Mostre que $G$ é abeliano.

        \item Seja $H$ um subgrupo de $G$. Tome $a \in G$ tal que $a \notin H$. Mostre que $H \cup aH$ é um subgrupo de $G$.

        \item Mostre que $|H \cup aH| = 2|H|$.
    \end{enumerate}

    \vspace{.3cm}

    \questao{} Considere o conjunto $\z_{24}$. Defina
    \[
        G = \{ \overline{a} \in \z_{24}^* \mid \mbox{ existe } \overline{b} \in \z_{24} \mbox{ tal que } \overline{a}\overline{b} = \overline{1}\}.
    \]
    \begin{enumerate}[label={\alph*})]
        \item  Mostre que $G$ é um grupo com a multiplicação de $\z_{24}$.

        \item Encontre todos os subgrupos cíclicos de $G$.
    \end{enumerate}

    \vspace{.3cm}

    \questao{} Considere o conjunto $\z_{20}$. Defina
    \[
        G = \{ \overline{a} \in \z_{20}^* \mid \mbox{ existe } \overline{b} \in \z_{20} \mbox{ tal que } \overline{a}\overline{b} = \overline{1}\}.
    \]
    \begin{enumerate}[label={\alph*})]
        \item  Mostre que $G$ é um grupo com a multiplicação de $\z_{20}$.

        \item Encontre dois subgrupos de ordem 4 em $G$. Algum deles é cíclico?
    \end{enumerate}
\end{document}
